\chapter{Lenses, Images, and the Eye}\label{ch:g11:lenses-and-eye}

Every sharp picture --- on a phone sensor, a cinema screen, or the back
of your own eye --- comes from the same trick: a curved transparent
medium bends rays by \omterm{def:g10:refraction:refraction}{refraction} (\cref{ch:g10:refraction}) so that all
the rays leaving one point meet again at another. This chapter tames
the trick: one \omterm{def:g11:lenses-and-eye:lens}{lens}, three rays, one formula --- then points it at
cameras, spectacles, and the living \omterm{def:g11:lenses-and-eye:lens}{lens} reading this sentence.

\section{Converging and diverging lenses}

\begin{definition}[Thin lens]\label{def:g11:lenses-and-eye:lens}
A \emph{lens}\index{lens} is a transparent medium bounded by two
surfaces, at least one curved. Thicker at the center than at the edge,
it is \emph{converging}\index{converging lens}: parallel rays are bent
toward each other; thicker at the edge, \emph{diverging}\index{diverging
lens}: they spread apart. A \emph{thin}\index{thin lens} lens has
negligible thickness: we draw a segment with outward (converging) or
inward (diverging) arrowheads. Its center $O$ is the \emph{optical
center}\index{optical center}; the perpendicular line through $O$, the
\emph{principal axis}\index{principal axis}.
\end{definition}

\begin{definition}[Focal points and focal length]\label{def:g11:lenses-and-eye:focal}
Rays arriving parallel to the axis of a \omterm{def:g11:lenses-and-eye:lens}{converging lens} all cross the
axis at one point behind it, the \emph{image focal point}\index{focal
point} $F'$; symmetrically, rays leaving the \emph{object focal point}
$F$ (in front, with $OF = OF'$) exit parallel to the axis. The
\emph{focal length}\index{focal length} $f' = \overline{OF'}$ is
positive for a \omterm{def:g11:lenses-and-eye:lens}{converging lens}; through a \omterm{def:g11:lenses-and-eye:lens}{diverging lens}, parallel rays
exit as if from a point $F'$ \emph{in front} of the \omterm{def:g11:lenses-and-eye:lens}{lens}: $f' < 0$.
\end{definition}

\begin{definition}[Vergence]\label{def:g11:lenses-and-eye:vergence}
The \emph{vergence}\index{vergence} of a \omterm{def:g11:lenses-and-eye:lens}{lens} of \omterm{def:g11:lenses-and-eye:focal}{focal length} $f'$
(in meters) is $C = 1/f'$,
measured in \emph{diopters}\index{diopter} (symbol D):
$1~\text{D} = \qty{1}{m^{-1}}$; \omterm{def:g11:lenses-and-eye:lens}{diverging} lenses have $C < 0$.
Opticians prescribe in diopters: vergences of \omterm{def:g11:lenses-and-eye:lens}{thin} lenses in contact
simply add (admitted here; derived in the Year 1 volume).
\end{definition}

\begin{example}[Reading a prescription]\label{ex:g11:lenses-and-eye:vergence}
A \omterm{def:g11:lenses-and-eye:lens}{lens} with $f' = \qty{0.50}{m}$ has $C = 1/0.50 = +2.0$~D:
\omterm{def:g11:lenses-and-eye:lens}{converging}. A $-4.0$~D prescription means $f' = \qty{-0.25}{m}$:
\omterm{def:g11:lenses-and-eye:lens}{diverging}, \omterm{def:g11:lenses-and-eye:focal}{focal length} \qty{25}{cm}.
\end{example}

\section{Constructing the image}

\begin{method}[The three construction rays]\label{met:g11:lenses-and-eye:rays}
Let $AB$ be an object perpendicular to the axis at $A$. Among all rays
leaving the tip $B$, three have known paths:
\begin{enumerate}
  \item the ray through the \omterm{def:g11:lenses-and-eye:lens}{optical center} $O$ goes straight on;
  \item the ray parallel to the axis exits through $F'$;
  \item the ray through $F$ exits parallel to the axis.
\end{enumerate}
Any two suffice: their intersection is the image point $B'$, and the
image $A'B'$ stands perpendicular to the axis at $A'$. If the outgoing
rays diverge, their \emph{backward} extensions (dashed) meet at a
\omterm{def:g11:lenses-and-eye:images}{virtual image} point.
\end{method}

\begin{omfigure}
\begin{tikzpicture}[scale=0.55]
  \draw[-Stealth] (-6.4,0) -- (6.4,0);
  \draw[Stealth-Stealth, omDef, very thick] (0,-2.4) -- (0,2.4);
  \fill (0,0) circle (1.5pt) node[below right, font=\small] {$O$};
  \fill (-2,0) circle (1.5pt) node[below left, font=\small] {$F$};
  \fill (2,0) circle (1.5pt) node[below right, font=\small] {$F'$};
  \draw[omMeth] (-4,1.2) -- (0,1.2) -- (4,-1.2);
  \draw[omMeth] (-4,1.2) -- (4,-1.2);
  \draw[omMeth] (-4,1.2) -- (0,-1.2) -- (4,-1.2);
  \draw[-Stealth, very thick] (-4,0) -- (-4,1.2);
  \node[above, font=\small] at (-4,1.2) {$B$};
  \node[below, font=\small] at (-4,0) {$A$};
  \draw[-Stealth, omThm, very thick] (4,0) -- (4,-1.2);
  \node[below, font=\small] at (4,-1.2) {$B'$};
  \node[above, font=\small] at (4,0) {$A'$};
\end{tikzpicture}

{\small The three construction rays: object $AB$ beyond $F$, image
$A'B'$ \omterm{def:g11:lenses-and-eye:images}{real} and inverted (here $\overline{OA} = -2f'$: life-size,
$\gamma = -1$).}
\end{omfigure}

\begin{definition}[Real and virtual images]\label{def:g11:lenses-and-eye:images}
An image is \emph{real}\index{real image} when the outgoing rays
actually pass through it: a screen placed there catches it;
\emph{virtual}\index{virtual image} when only their backward extensions
meet: no screen catches it, but an eye looking into the \omterm{def:g11:lenses-and-eye:lens}{lens} sees it
perfectly well. A \omterm{def:g11:lenses-and-eye:lens}{converging lens} gives a real, inverted image of any
object beyond $F$; slide the object inside the \omterm{def:g11:lenses-and-eye:focal}{focal length} and the
image turns virtual, upright, enlarged.
\end{definition}

\begin{omfigure}
\begin{tikzpicture}[scale=0.5]
  \draw[-Stealth] (-7.2,0) -- (4.6,0);
  \draw[Stealth-Stealth, omDef, very thick] (0,-2.2) -- (0,3.6);
  \fill (0,0) circle (1.5pt) node[below right, font=\small] {$O$};
  \fill (-2,0) circle (1.5pt) node[below, font=\small] {$F$};
  \fill (2,0) circle (1.5pt) node[below, font=\small] {$F'$};
  \draw[omMeth] (-1.5,0.8) -- (3,-1.6);
  \draw[omMeth] (-1.5,0.8) -- (0,0.8) -- (4,-0.8);
  \draw[omMeth, dashed] (0,0.8) -- (-6,3.2);
  \draw[omMeth, dashed] (-1.5,0.8) -- (-6,3.2);
  \draw[-Stealth, very thick] (-1.5,0) -- (-1.5,0.8);
  \node[below, font=\small] at (-1.5,0) {$A$};
  \node[above right, font=\small] at (-1.4,0.75) {$B$};
  \draw[-Stealth, omThm, very thick, dashed] (-6,0) -- (-6,3.2);
  \node[below, font=\small] at (-6,0) {$A'$};
  \node[above, font=\small] at (-6,3.2) {$B'$};
\end{tikzpicture}

{\small Object inside the \omterm{def:g11:lenses-and-eye:focal}{focal length} (the \omterm{def:g11:lenses-and-eye:magnifier}{magnifier} setting): the
backward extensions build a \omterm{def:g11:lenses-and-eye:images}{virtual}, upright, enlarged image $A'B'$.}
\end{omfigure}

\section{One formula: the conjugate relation}

Constructions show \emph{where} the image is; a formula computes it.
Orient the axis along the light and use \emph{algebraic measures}:
$\overline{OA}$ is the coordinate of $A$ relative to $O$ (negative for
a \omterm{def:g11:lenses-and-eye:images}{real} object in front); heights $\overline{AB}$ count positive upward.

\begin{theorem}[Thin-lens conjugate relation]\label{thm:g11:lenses-and-eye:conjugate}
Through a \omterm{def:g11:lenses-and-eye:lens}{thin lens} of \omterm{def:g11:lenses-and-eye:focal}{focal length} $f'$, an object point $A$ on the
axis images at the point $A'$ with
\[
\frac{1}{\overline{OA'}} - \frac{1}{\overline{OA}} = \frac{1}{f'} = C .
\]
\end{theorem}
\admitted

\begin{remark}
The relation follows from \omterm{thm:g10:refraction:snell}{Snell's law} applied to rays close to the
axis; the derivation --- and the limits of the \omterm{def:g11:lenses-and-eye:lens}{thin-lens} model --- are
made honest in the Year 1 volume.
\end{remark}

\begin{definition}[Magnification]\label{def:g11:lenses-and-eye:magnification}
The \emph{magnification}\index{magnification} of the image is
\[
\gamma = \frac{\overline{A'B'}}{\overline{AB}}
       = \frac{\overline{OA'}}{\overline{OA}} .
\]
$\abs{\gamma}$ compares sizes, its sign orientations ($\gamma < 0$:
inverted); the sign of $\overline{OA'}$ reads the nature (positive:
\omterm{def:g11:lenses-and-eye:images}{real}, behind the \omterm{def:g11:lenses-and-eye:lens}{lens}; negative: \omterm{def:g11:lenses-and-eye:images}{virtual}, in front).
\end{definition}

\begin{example}[The formula at work]\label{ex:g11:lenses-and-eye:conjugate}
An object stands \qty{30}{cm} in front of a \omterm{def:g11:lenses-and-eye:lens}{converging lens} with
$f' = \qty{10}{cm}$: $\overline{OA} = \qty{-30}{cm}$, so
\[
\frac{1}{\overline{OA'}} = \frac{1}{10} - \frac{1}{30} = \frac{1}{15},
\qquad \overline{OA'} = +\qty{15}{cm},
\qquad \gamma = \frac{15}{-30} = -0.50 :
\]
a \omterm{def:g11:lenses-and-eye:images}{real image} \qty{15}{cm} behind the \omterm{def:g11:lenses-and-eye:lens}{lens}, inverted, half-size ---
exactly what the three rays draw.
\end{example}

\section{The eye}

\begin{definition}[The eye as an optical system]\label{def:g11:lenses-and-eye:eye}
Optically, the eye is a \omterm{def:g11:lenses-and-eye:lens}{converging lens} facing a screen: the
\emph{cornea}\index{cornea} and \emph{crystalline lens}\index{crystalline
lens} act as one \omterm{def:g11:lenses-and-eye:lens}{converging lens} of adjustable \omterm{def:g11:lenses-and-eye:focal}{focal length}; the
\emph{retina}\index{retina}, about \qty{17}{mm} behind, is the fixed
screen. Focusing is done by squeezing the \omterm{def:g11:lenses-and-eye:lens}{lens}, not by moving it as in
a camera: \emph{accommodation}\index{accommodation} is the change of
$f'$ produced by the ciliary muscles. The closest point seen sharply at full
accommodation is the \emph{near point}\index{near point} --- about
\qty{25}{cm} for the \emph{standard eye}\index{standard eye} --- the
farthest, at rest, the \emph{far point}\index{far point} (infinity for
the standard eye).
\end{definition}

\begin{omfigure}
\begin{tikzpicture}[scale=0.85]
  \draw[thick] (0,0) circle (1.1);
  \draw[omDef, very thick] (-0.95,-0.5) -- (-0.95,0.5);
  \draw (-2.6,0.35) -- (-0.95,0.35) -- (1.1,0);
  \draw (-2.6,0) -- (1.1,0);
  \draw (-2.6,-0.35) -- (-0.95,-0.35) -- (1.1,0);
  \fill (1.1,0) circle (1.5pt) node[right, font=\small] {retina};
  \node[font=\small] at (-0.7,-1.5) {at rest: distant object};
\end{tikzpicture}\qquad
\begin{tikzpicture}[scale=0.85]
  \draw[thick] (0,0) circle (1.1);
  \draw[omDef, line width=2.6pt] (-0.95,-0.5) -- (-0.95,0.5);
  \fill (-2.6,0) circle (1.5pt) node[above, font=\small] {$B$};
  \draw (-2.6,0) -- (-0.95,0.4) -- (1.1,0);
  \draw (-2.6,0) -- (1.1,0);
  \draw (-2.6,0) -- (-0.95,-0.4) -- (1.1,0);
  \node[font=\small] at (-0.7,-1.5) {accommodated: near object};
\end{tikzpicture}

{\small \omterm{def:g11:lenses-and-eye:eye}{Accommodation}: to focus a near object on the same \omterm{def:g11:lenses-and-eye:eye}{retina}, the
\omterm{def:g11:lenses-and-eye:eye}{crystalline lens} bulges --- its \omterm{def:g11:lenses-and-eye:vergence}{vergence} increases.}
\end{omfigure}

\begin{example}[The vergence of your eye]\label{ex:g11:lenses-and-eye:eyerange}
The \omterm{def:g11:lenses-and-eye:eye}{retina} sits at $\overline{OA'} = +\qty{0.017}{m}$. Distant object
($1/\overline{OA} \approx 0$): $C = 1/0.017 \approx 59$~D. \omterm{def:g11:lenses-and-eye:eye}{Near point}
($\overline{OA} = \qty{-0.25}{m}$): $C = 1/0.017 + 1/0.25 \approx
63$~D. Reading costs about $4$~D: the eye is a $59$~D \omterm{def:g11:lenses-and-eye:lens}{lens} with a
built-in $+4$~D fine tune.
\end{example}

\section{Defects, corrections, and the magnifier}

\begin{definition}[Myopia and hyperopia]\label{def:g11:lenses-and-eye:defects}
A \emph{myopic}\index{myopia} (short-sighted) eye is too \omterm{def:g11:lenses-and-eye:lens}{converging} for
its length: parallel rays focus \emph{in front of} the \omterm{def:g11:lenses-and-eye:eye}{retina}, distant
objects blur, the \omterm{def:g11:lenses-and-eye:eye}{far point} is at finite distance; a \omterm{def:g11:lenses-and-eye:lens}{diverging}
spectacle \omterm{def:g11:lenses-and-eye:lens}{lens} corrects it. A \emph{hyperopic}\index{hyperopia}
(far-sighted) eye is not \omterm{def:g11:lenses-and-eye:lens}{converging} enough: near objects would focus
behind the \omterm{def:g11:lenses-and-eye:eye}{retina}, the \omterm{def:g11:lenses-and-eye:eye}{near point} recedes; a \omterm{def:g11:lenses-and-eye:lens}{converging lens} corrects
it. Aging stiffens the \omterm{def:g11:lenses-and-eye:eye}{crystalline lens}
(\emph{presbyopia}\index{presbyopia}), shrinking \omterm{def:g11:lenses-and-eye:eye}{accommodation} --- same
remedy: reading glasses.
\end{definition}

\begin{example}[Prescribing for a myope]\label{ex:g11:lenses-and-eye:correction}
A \omterm{def:g11:lenses-and-eye:defects}{myopic} eye has its \omterm{def:g11:lenses-and-eye:eye}{far point} at \qty{50}{cm}: the correcting \omterm{def:g11:lenses-and-eye:lens}{lens}
must show distant objects \emph{at the \omterm{def:g11:lenses-and-eye:eye}{far point}} --- object at
infinity, image at $\overline{OA'} = \qty{-0.50}{m}$ (\omterm{def:g11:lenses-and-eye:images}{virtual}, in
front). Then $1/f' = 1/\overline{OA'}$: $f' = \qty{-0.50}{m}$,
$C = -2.0$~D; the relaxed eye looks at that \omterm{def:g11:lenses-and-eye:images}{virtual image} and sees it
sharply.
\end{example}

\begin{omfigure}
\begin{tikzpicture}[scale=0.85]
  \draw[thick] (0,0) circle (1.1);
  \draw[omDef, very thick] (-0.95,-0.5) -- (-0.95,0.5);
  \draw (-2.7,0.35) -- (-0.95,0.35) -- (0.55,0) -- (1.1,-0.13);
  \draw (-2.7,0) -- (1.1,0);
  \draw (-2.7,-0.35) -- (-0.95,-0.35) -- (0.55,0) -- (1.1,0.13);
  \fill[omThm] (0.55,0) circle (1.5pt);
  \node[font=\small] at (-0.7,-1.5) {myopic: focus before the retina};
\end{tikzpicture}\qquad
\begin{tikzpicture}[scale=0.85]
  \draw[thick] (0,0) circle (1.1);
  \draw[omDef, very thick] (-0.95,-0.5) -- (-0.95,0.5);
  \draw[omMeth, very thick, {Stealth[reversed]}-{Stealth[reversed]}]
    (-1.9,-0.55) -- (-1.9,0.55);
  \draw (-2.8,0.3) -- (-1.9,0.3) -- (-0.95,0.42) -- (1.1,0);
  \draw (-2.8,0) -- (1.1,0);
  \draw (-2.8,-0.3) -- (-1.9,-0.3) -- (-0.95,-0.42) -- (1.1,0);
  \node[font=\small] at (-0.7,-1.5) {corrected: diverging lens};
\end{tikzpicture}

{\small A \omterm{def:g11:lenses-and-eye:defects}{myopic} eye focuses parallel rays before the \omterm{def:g11:lenses-and-eye:eye}{retina} (red
point); a \omterm{def:g11:lenses-and-eye:lens}{diverging lens} lands the focus on the \omterm{def:g11:lenses-and-eye:eye}{retina}.}
\end{omfigure}

\begin{definition}[The magnifier]\label{def:g11:lenses-and-eye:magnifier}
A \emph{magnifier}\index{magnifier} is a \omterm{def:g11:lenses-and-eye:lens}{converging lens} of short \omterm{def:g11:lenses-and-eye:focal}{focal
length} used with the object \emph{inside} the \omterm{def:g11:lenses-and-eye:focal}{focal length}: the image
is \omterm{def:g11:lenses-and-eye:images}{virtual}, upright, enlarged (\cref{def:g11:lenses-and-eye:images}),
and the eye views it comfortably far away.
\end{definition}

\begin{example}[A jeweler's loupe]\label{ex:g11:lenses-and-eye:magnifier}
Loupe with $f' = \qty{5.0}{cm}$, stone at \qty{4.0}{cm}:
$1/\overline{OA'} = 1/5.0 - 1/4.0 = -1/20$, so
$\overline{OA'} = \qty{-20}{cm}$ and $\gamma = (-20)/(-4.0) = +5.0$: a
\omterm{def:g11:lenses-and-eye:images}{virtual}, upright image, five times larger, a comfortable \qty{20}{cm}
from the \omterm{def:g11:lenses-and-eye:lens}{lens}.
\end{example}

\section{Exercises}

\begin{exercise}[$\star$]\label{exo:g11:lenses-and-eye:1}
Compute the \omterm{def:g11:lenses-and-eye:vergence}{vergences} for $f' = \qty{50}{cm}$, $\qty{-20}{cm}$ and
$\qty{4.0}{mm}$, then the \omterm{def:g11:lenses-and-eye:focal}{focal length} of a $+8.0$~D \omterm{def:g11:lenses-and-eye:lens}{lens}.
\end{exercise}

\begin{exercise}[$\star$]\label{exo:g11:lenses-and-eye:2}
A \omterm{def:g11:lenses-and-eye:lens}{lens} is thicker at the edges than at the center. \omterm{def:g11:lenses-and-eye:lens}{Converging} or
\omterm{def:g11:lenses-and-eye:lens}{diverging}? Can it serve as a \omterm{def:g11:lenses-and-eye:magnifier}{magnifier}? What do you see holding it
over a printed page?
\end{exercise}

\begin{exercise}[$\star$]\label{exo:g11:lenses-and-eye:3}
A \omterm{def:g11:lenses-and-eye:lens}{converging lens} has $f' = \qty{10}{cm}$; an object stands
\qty{20}{cm} in front of it. Construct the image with two of the three
rays, then check its position and size with the conjugate relation.
\end{exercise}

\begin{exercise}[$\star$]\label{exo:g11:lenses-and-eye:4}
An object stands \qty{24}{cm} in front of a \omterm{def:g11:lenses-and-eye:lens}{converging lens} with
$f' = \qty{8.0}{cm}$. Compute $\overline{OA'}$ and $\gamma$, and
describe the image (nature, orientation, size).
\end{exercise}

\begin{exercise}[$\star$]\label{exo:g11:lenses-and-eye:5}
In the eye: what forms the \omterm{def:g11:lenses-and-eye:lens}{converging lens}? What plays the screen? What
changes during \omterm{def:g11:lenses-and-eye:eye}{accommodation} --- and what cannot? Define the \omterm{def:g11:lenses-and-eye:eye}{near point}
of the \omterm{def:g11:lenses-and-eye:eye}{standard eye}.
\end{exercise}

\begin{exercise}[$\star\star$]\label{exo:g11:lenses-and-eye:6}
A camera \omterm{def:g11:lenses-and-eye:lens}{lens} has $f' = \qty{50}{mm}$. Where is the sensor when focused
on a distant landscape? Compute the lens--sensor distance for a subject
\qty{3.0}{m} away, and the travel between the two settings.
\end{exercise}

\begin{exercise}[$\star\star$]\label{exo:g11:lenses-and-eye:7}
A \omterm{def:g11:lenses-and-eye:magnifier}{magnifier} has $f' = \qty{6.0}{cm}$; a stamp lies \qty{5.0}{cm} below
it. Position, nature, orientation and \omterm{def:g11:lenses-and-eye:magnification}{magnification} of the image?
\end{exercise}

\begin{exercise}[$\star\star$]\label{exo:g11:lenses-and-eye:8}
An object and a screen face each other; a \omterm{def:g11:lenses-and-eye:lens}{converging lens} between them
gives a sharp image when both are \qty{40}{cm} from the \omterm{def:g11:lenses-and-eye:lens}{lens}. Compute
$f'$ and $\gamma$. Why must this symmetric image be exactly life-size?
\end{exercise}

\begin{exercise}[$\star\star$]\label{exo:g11:lenses-and-eye:9}
An object stands \qty{50}{cm} in front of a \omterm{def:g11:lenses-and-eye:lens}{diverging lens} with
$f' = \qty{-25}{cm}$. Compute $\overline{OA'}$ and $\gamma$; describe
the image. Why do a myope's eyes look slightly smaller through their
glasses?
\end{exercise}

\begin{exercise}[$\star\star$]\label{exo:g11:lenses-and-eye:10}
A \omterm{def:g11:lenses-and-eye:defects}{myopic} eye has its \omterm{def:g11:lenses-and-eye:eye}{far point} at \qty{40}{cm}. What \omterm{def:g11:lenses-and-eye:vergence}{vergence} must the
correcting \omterm{def:g11:lenses-and-eye:lens}{lens} have, and where does it put the image of a distant
mountain? Why is the mountain then seen effortlessly?
\end{exercise}

\begin{exercise}[$\star\star$]\label{exo:g11:lenses-and-eye:11}
A presbyopic reader has their \omterm{def:g11:lenses-and-eye:eye}{near point} at \qty{80}{cm} but wants to
read at \qty{25}{cm}: the glasses must image the page at \qty{25}{cm}
onto a \omterm{def:g11:lenses-and-eye:images}{virtual} page at \qty{80}{cm}. Compute the required \omterm{def:g11:lenses-and-eye:vergence}{vergence}.
\end{exercise}

\begin{exercise}[$\star\star\star$]\label{exo:g11:lenses-and-eye:12}
A projector must throw an image of a slide, magnified $\gamma = -50$,
onto a screen \qty{5.1}{m} from the slide.
\begin{enumerate}
  \item Find $\overline{OA}$, $\overline{OA'}$, given
        $\overline{OA'} - \overline{OA} = \qty{5.1}{m}$.
  \item Deduce the \omterm{def:g11:lenses-and-eye:focal}{focal length} of the projection \omterm{def:g11:lenses-and-eye:lens}{lens}.
  \item How tall is the image of a \qty{24}{mm}-tall slide?
\end{enumerate}
\end{exercise}

\begin{exercise}[$\star\star\star$]\label{exo:g11:lenses-and-eye:13}
The \omterm{def:g11:lenses-and-eye:eye}{retina} is \qty{17}{mm} behind the eye's \omterm{def:g11:lenses-and-eye:lens}{lens}.
\begin{enumerate}
  \item Compute the eye's \omterm{def:g11:lenses-and-eye:vergence}{vergence} focused at infinity, then on the
        standard \omterm{def:g11:lenses-and-eye:eye}{near point} (\qty{25}{cm}). Deduce the \omterm{def:g11:lenses-and-eye:eye}{accommodation}
        \omterm{def:g10:signals-and-waves:amplitude}{amplitude} of the \omterm{def:g11:lenses-and-eye:eye}{standard eye}.
  \item An older eye has its \omterm{def:g11:lenses-and-eye:eye}{near point} at \qty{1.0}{m}. Compute its
        \omterm{def:g10:signals-and-waves:amplitude}{amplitude} and the fraction of the standard \omterm{def:g10:signals-and-waves:amplitude}{amplitude} remaining.
\end{enumerate}
\end{exercise}

\begin{exercise}[$\star\star\star$]\label{exo:g11:lenses-and-eye:14}
A \omterm{def:g11:lenses-and-eye:magnifier}{magnifier} sold as ``$\times 3$'' has its \omterm{def:g11:lenses-and-eye:focal}{focal length} given by
$3 = 0.25/f'$ ($f'$ in meters).
\begin{enumerate}
  \item Compute $f'$.
  \item Where must an object sit for its \omterm{def:g11:lenses-and-eye:images}{virtual image} to lie
        \qty{25}{cm} in front of the \omterm{def:g11:lenses-and-eye:lens}{lens}? Compute $\gamma$; compare
        with the advertised $\times 3$.
\end{enumerate}
\end{exercise}

\begin{exercise}[$\star\star\star$]\label{exo:g11:lenses-and-eye:15}
A \omterm{def:g11:lenses-and-eye:images}{real} object stands at distance $d > 0$ in front of a \omterm{def:g11:lenses-and-eye:lens}{converging lens}
($\overline{OA} = -d$, $f' > 0$). Show that
$\overline{OA'} = \frac{f'd}{d - f'}$ and
$\gamma = \frac{f'}{f' - d}$. Deduce: \omterm{def:g11:lenses-and-eye:images}{real} and inverted when $d > f'$;
\omterm{def:g11:lenses-and-eye:images}{virtual}, upright, \emph{enlarged} ($\gamma > 1$) when $0 < d < f'$ ---
every \omterm{def:g11:lenses-and-eye:lens}{converging lens} is a \omterm{def:g11:lenses-and-eye:magnifier}{magnifier}, used close enough.
\end{exercise}

\section{Problem: One formula, three instruments}

\begin{problem}[{Weekend problem --- the same conjugate relation
focuses a camera, prescribes reading glasses, designs a loupe, and
measures the zoom built into your own eye}]
\label{pb:g11:lenses-and-eye:1}
A photographer racks a \qty{50}{mm} \omterm{def:g11:lenses-and-eye:lens}{lens}, a grandmother holds the
newspaper at arm's length, a jeweler leans over a diamond: one line of
algebra, $1/\overline{OA'} - 1/\overline{OA} = 1/f'$
(\cref{thm:g11:lenses-and-eye:conjugate}), runs all three trades ---
and then the eye itself.

\medskip
\noindent\textbf{Part I --- One formula.} A \omterm{def:g11:lenses-and-eye:lens}{converging lens} has
$f' = \qty{20}{cm}$.
\begin{enumerate}
  \item An object at \qty{60}{cm}: compute $\overline{OA'}$ and
        $\gamma$; describe the image.
  \item Move the object to \qty{30}{cm}; recompute. What got exchanged
        with question 1? (Light paths are reversible.)
  \item Move it to \qty{10}{cm}, inside $f'$: recompute and describe.
  \item Summarize: as the object slides in from very far toward $F$,
        then past it, how does the image move and change nature?
  \item Show that a very distant object ($1/\overline{OA} \approx 0$)
        always images in the focal plane: $\overline{OA'} = f'$.
\end{enumerate}

\medskip
\noindent\textbf{Part II --- The photographer.}
A camera \omterm{def:g11:lenses-and-eye:lens}{lens} has $f' = \qty{50}{mm}$; the sensor, \qty{24}{mm} tall,
can be moved relative to the \omterm{def:g11:lenses-and-eye:lens}{lens}.
\begin{enumerate}[resume]
  \item Where is the sensor when focused on a distant landscape?
  \item Now a face \qty{1.0}{m} away: compute the new lens--sensor
        distance and the travel from the landscape setting.
  \item The mechanism allows at most \qty{55}{mm} of lens--sensor
        distance. Compute the minimum focusing distance.
  \item At that closest focus, compute $\gamma$. Does the image of a
        \qty{20}{cm} face fit on the sensor?
  \item Life-size ``macro'' means $\gamma = -1$: where must object and
        sensor then be, and why does this demand a special \omterm{def:g11:lenses-and-eye:lens}{lens}?
\end{enumerate}

\medskip
\noindent\textbf{Part III --- The grandmother.}
With age the \omterm{def:g11:lenses-and-eye:eye}{near point} recedes: grandmother's is at \qty{1.0}{m}.
\begin{enumerate}[resume]
  \item Explain the arm's-length newspaper.
  \item To read at \qty{25}{cm}, her glasses must turn a page at
        \qty{25}{cm} into a \omterm{def:g11:lenses-and-eye:images}{virtual image} at \qty{1.0}{m}: compute the
        required \omterm{def:g11:lenses-and-eye:vergence}{vergence}.
  \item Compute $\gamma$. The image is four times larger, yet looks no
        bigger: explain, comparing sizes \emph{and} distances.
  \item Show that print is sharp only between \qty{25}{cm} and
        $f' \approx \qty{33}{cm}$. What of a \emph{distant} object seen
        through the glasses?
  \item Deduce in one sentence why bifocal lenses exist.
\end{enumerate}

\medskip
\noindent\textbf{Part IV --- The jeweler, and the eye's own zoom.}
The jeweler's loupe has $f' = \qty{5.0}{cm}$; his \omterm{def:g11:lenses-and-eye:eye}{near point} is at
\qty{25}{cm}; his \omterm{def:g11:lenses-and-eye:eye}{retina} is \qty{17}{mm} behind his eye's \omterm{def:g11:lenses-and-eye:lens}{lens}.
\begin{enumerate}[resume]
  \item A stone \qty{4.0}{cm} under the loupe: position, nature,
        \omterm{def:g11:lenses-and-eye:magnification}{magnification}?
  \item Where must the stone sit for its image to fall at his \omterm{def:g11:lenses-and-eye:eye}{near
        point}, and what is $\gamma$ then?
  \item Why do experienced jewelers place the stone exactly at $F$
        instead (image at infinity)?
  \item Now the eye alone: compute its \omterm{def:g11:lenses-and-eye:vergence}{vergence} viewing a distant
        object, then one at \qty{25}{cm}: how many \omterm{def:g11:lenses-and-eye:vergence}{diopters} of
        \omterm{def:g11:lenses-and-eye:eye}{accommodation}?
  \item Convert both \omterm{def:g11:lenses-and-eye:vergence}{vergences} to \omterm{def:g11:lenses-and-eye:focal}{focal lengths}: by how much, in
        millimeters and percent, does the eye change its own \omterm{def:g11:lenses-and-eye:focal}{focal
        length} --- the zoom the camera did with travel and grandmother
        with $+3$~D?
\end{enumerate}
\end{problem}
